\section*{Kurzfassung}
\addcontentsline{toc}{section}{Kurzfassung}

Bei der Transformation heterogener Engineering-Bestandsinformation in SysML v2
müssen unterschiedlich strukturierte Systembeschreibungen in eine formale
Repräsentation überführt werden. Menschliches Expertenwissen bleibt dabei
insbesondere für fachliche und modellbezogene Entscheidungen wesentlich.
KI-basierte probabilistische Verfahren können jedoch Interpretationen bestehender
Informationen erzeugen und damit Entscheidungsgrundlagen für nachgelagerte
menschliche Bewertungen bereitstellen.

Diese Arbeit entwickelt und untersucht eine Informationsverarbeitungsarchitektur
für einen kontrollierten KI-gestützten Transformationsprozess. Eine definierte
Informationsbasis unterscheidet Engineering- und Kontextquellen und bindet
abgeleitete Aussagen über quellengebundene Evidenz an ihre Herkunft.
Mehrperspektivische KI-Interpretationen dienen als Entscheidungsgrundlage für
menschliche Experten. Die Human Engineering Authority autorisiert die fachliche
Verwendung einer Information, während die Model Authority über deren strukturelle
Einordnung und Repräsentation entscheidet. Autorisierte Inhalte werden in einem
Internal Engineering Model zusammengeführt und anschließend regelgebunden in
SysML v2 überführt. Provenienz und Nachverfolgbarkeit verbinden dabei Quellen,
Interpretationen, menschliche Entscheidungen und Modellzustände über den gesamten
Transformationsprozess.

Das Architekturkonzept wurde prototypisch realisiert und durch wiederholte
Transformationsläufe untersucht. Dabei konnten Informations-, Entscheidungs- und
Modellbildungsfunktionen als eigenständige Architekturgrenzen identifiziert und
präzisiert werden. Die Einbeziehung mehrerer Engineering-Quellen erweiterte die
Untersuchung um den Übergang von quellenbezogen autorisierten Informationszuständen
zu einer gemeinsamen projektweiten Modellbildung. Ein zusammenhängender
Verarbeitungspfad von registrierten Engineering-Quellen über KI-gestützte
Interpretation und menschliche Autorisierung bis zu einem technisch validierten
und menschlich freigegebenen SysML-v2-Artefakt konnte ausgeführt und anhand
persistierter Zustände rekonstruiert werden.

Die Ergebnisse zeigen, dass probabilistische KI-Verfahren für die Transformation
heterogener Engineering-Bestandsinformation nutzbar gemacht werden können, wenn
ihre Ergebnisse in einen kontrollierten Rahmen aus expliziter menschlicher
Autorisierung, semantischen und modellbezogenen Leitplanken, regelgebundener
Transformation sowie technischer Validierung eingebettet werden.

\clearpage
\section*{Abstract}
\addcontentsline{toc}{section}{Abstract}

The transformation of heterogeneous legacy engineering information into SysML v2
requires differently structured system descriptions to be transferred into a formal
representation. Human expertise remains essential, particularly for domain-specific
and model-related decisions. However, AI-based probabilistic methods can generate
interpretations of existing information and thereby provide decision support for
subsequent human evaluation.

This thesis develops and investigates an information processing architecture for a
controlled AI-supported transformation process. A defined information basis
distinguishes between engineering sources and contextual sources and links derived
statements to their origin through source-grounded evidence. Multi-perspective AI
interpretations provide decision support for human experts. The Human Engineering
Authority authorizes the domain-specific use of information, while the Model Authority
decides on its structural placement and representation within the model. Authorized
content is consolidated in an Internal Engineering Model and subsequently transformed
into SysML v2 in a rule-based manner. Provenance and traceability connect sources,
interpretations, human decisions, and model states throughout the entire transformation
process.

The architecture was implemented as a prototype and investigated through repeated
transformation runs. This made it possible to identify and refine information-processing,
decision-making, and model-building functions as distinct architectural boundaries.
The inclusion of multiple engineering sources extended the investigation to the
transition from source-specific authorized information states to joint project-level
model formation. A continuous processing path from registered engineering sources
through AI-supported interpretation and human authorization to a technically validated
and human-approved SysML v2 artifact was executed and reconstructed on the basis of
persisted states.

The results show that probabilistic AI methods can be applied to the transformation of
heterogeneous legacy engineering information when their outputs are embedded in a
controlled framework combining explicit human authorization, semantic and
model-related constraints, rule-based transformation, and technical validation.